% =========================================================================
%  DE RERUM TODO EXISTENS — RESTRUCTURED EDITION
%  Shared Preamble
% =========================================================================

\usepackage{amsmath,amssymb,amsfonts}
\usepackage{mathtools}
\usepackage{booktabs}
\usepackage{longtable}
\usepackage{graphicx}
\usepackage[table,dvipsnames]{xcolor}
\usepackage{hyperref}
\usepackage{cleveref}
\usepackage{fancyhdr}
\usepackage{titlesec}
\usepackage[margin=2.5cm]{geometry}
\usepackage{pifont}   % for \ding checkmark

% =========================================================================
%  NOTATION: THE KOPPA DESCRIPTOR FAMILY
% =========================================================================
%
%  Three related but DISTINCT dimensionless quantities:
%
%  1.  K  ≡  c/v        — the KINEMATIC RATIO of any specific orbit
%                          (a state variable; different for every orbit)
%
%  2.  K_body            — the kinematic ratio at a body's characteristic
%                          surface velocity (a system parameter):
%                            K_☉  = 686.5   (Sun)
%                            K_J  = 7,124   (Jupiter)
%                            K_⊕  = 37,848  (Earth)
%
%  3.  ϟ  ≡  (1/α)√(R_p/a_0) = 0.5464
%                        — the UNIVERSAL ATOMIC BRIDGE CONSTANT
%                          (a fundamental constant composed of R_p, a_0, α)
%
%  K is a measurement.  K_body is a property.  ϟ is a constant.
%  They are cousins, not reincarnations.
% =========================================================================

% --- Koppa symbol (the universal constant ϟ = 0.5464) ---
% U+03DF (ϟ) — using varkappa as closest standard LaTeX glyph
\newcommand{\kop}{\varkappa}
\newcommand{\kopfull}{koppa}

% --- Kinematic ratio for specific bodies ---
% Use these for body-specific values, NOT \kop
\newcommand{\Ksun}{K_\odot}
\newcommand{\Kjup}{K_J}
\newcommand{\Ksat}{K_S}
\newcommand{\Kearth}{K_\oplus}

% --- Physical constants ---
\newcommand{\Rp}{R_p}              % proton charge radius
\newcommand{\abohr}{a_0}           % Bohr radius
\newcommand{\fsc}{\alpha}          % fine structure constant
\newcommand{\Zeff}{Z_{\text{eff}}} % effective nuclear charge

% --- Formatting ---
\newcommand{\SDT}{\textsc{sdt}}
\renewcommand{\checkmark}{\ding{51}}

% --- Chapter styling ---
\titleformat{\chapter}[display]
  {\normalfont\huge\bfseries}
  {\chaptertitlename\ \thechapter}{20pt}{\Huge}

% --- Headers ---
\pagestyle{fancy}
\fancyhf{}
\fancyhead[LE,RO]{\thepage}
\fancyhead[LO]{\nouppercase{\rightmark}}
\fancyhead[RE]{\nouppercase{\leftmark}}

% --- Hyperref setup ---
\hypersetup{
  colorlinks=true,
  linkcolor=blue!70!black,
  citecolor=green!50!black,
  urlcolor=blue!80!black,
  pdftitle={De Rerum Todo Existens — Restructured Edition},
  pdfauthor={James Tyndall},
}

